% Generated draft: requires mathematical and source review.
\documentclass[11pt]{article}
\usepackage[margin=25mm]{geometry}
\usepackage{fontspec}
\setmainfont{Latin Modern Roman}
\usepackage{amsmath,amssymb,mathtools,unicode-math}
\setmathfont{Latin Modern Math}
\usepackage{enumitem,needspace,xurl,hyperref}
\setlength{\parindent}{0pt}
\setlength{\parskip}{5pt}
\setlength{\emergencystretch}{3em}
\newcommand{\N}{\mathbb{N}}
\newcommand{\Z}{\mathbb{Z}}
\newcommand{\Q}{\mathbb{Q}}
\newcommand{\R}{\mathbb{R}}
\newcommand{\C}{\mathbb{C}}
\newcommand{\F}{\mathbb{F}}
\newcommand{\A}{\mathbb{A}}
\newcommand{\PP}{\mathbb{P}}
\newcommand{\E}{\mathbb{E}}
\newcommand{\eps}{\varepsilon}
\newcommand{\dd}{\,\mathrm{d}}
\begin{document}
\section*{UnsolvedMath ID: 30004860}
\textbf{Source problem code:} catalogue entry
\subsection{Definitions and mathematical statement}
Let \(A\) be a \(*\)-algebra, with involution denoted by \(a\mapsto a^*\). The source defines positivity through all \(*\)-representations on Hilbert spaces, but does not specify the scalar field, whether representations must be bounded, or conventions concerning units.

Under the conventional bounded complex interpretation, a \(*\)-representation is a complex-linear multiplicative map
\[
\pi:A\longrightarrow\mathcal B(H),\qquad \pi(a^*)=\pi(a)^*,
\]
where \(H\) is a complex Hilbert space and \(\mathcal B(H)\) is its algebra of bounded linear operators. A bounded operator \(T\) is positive if \(T=T^*\) and \(\langle T\xi,\xi\rangle\geq 0\) for every \(\xi\in H\). With this interpretation, define
\[
A_{\mathrm{rep},+}
=
\left\{a\in A:\ \text{for every complex Hilbert space }H
\text{ and every }\pi:A\to\mathcal B(H)\text{ as above},\ 
\pi(a)=\pi(a)^*\text{ and }
\langle\pi(a)\xi,\xi\rangle\geq0\text{ for every }\xi\in H
\right\}.
\]
This interpretation is conditional: the supplied record does not establish that these are precisely the representation conventions intended in the report. If unbounded representations are intended, their domains and the meaning of positivity must first be specified.

\textbf{Problem.} How can one test whether a given element \(a\) of a \(*\)-algebra \(A\) maps to a positive operator in every \(*\)-representation on a Hilbert space? Under the bounded complex interpretation above, the target is to determine whether
\[
a\in A_{\mathrm{rep},+}.
\]
To make this an algorithmic decision problem, one must additionally specify a class of algebras, effective descriptions of the algebras and their elements, and the required form of a terminating decision procedure. These data are absent from the supplied formulation, so it does not specify a single algorithmic decision problem for arbitrary \(*\)-algebras. No assumption that \(a=a^*\) in \(A\) is supplied; the positivity requirement is imposed on its images.

The supplied dated curation labels the problem partially solved and reports effective procedures for free-group positivity, together with coRE-hardness and algorithmic undecidability for positivity in tensor products of free algebras. These are supplied literature assessments, not independently verified assertions here; the precise classes, encodings, and positivity conventions of those special cases are not given in the record.
\subsection{Short English statement}
The question asks how to recognize elements of a \(*\)-algebra whose images are positive operators in every Hilbert-space \(*\)-representation. A precise decision problem also needs representation conventions and effective input descriptions, which the supplied record does not specify. Its partially solved label and special-case complexity claims are supplied dated curation requiring review.
\subsection{Sources}
\begin{itemize}
\item Supplied local UnsolvedMath registry record, permanent ID \(30004860\), title \emph{Deciding Positivity in Star Algebras}, source problem code \texttt{catalogue entry}; collection: \emph{Oberwolfach Reports Open Problems}. The record lists proposed year 2021 and status \texttt{partially\_solved}.
\item Supplied report citation: \emph{Geometry and Optimization in Quantum Information} (2022), Oberwolfach Reports, pp.\ 2665--2720, DOI 10.4171/owr/2021/49. \url{https://doi.org/10.4171/owr/2021/49}. The background describes Oberwolfach Reports as documenting mathematical workshops and their open-problem discussions.
\item Supplied talk attribution: William Slofstra, \emph{Decision problems for positivity and sums of squares}, joint work with Arthur Mehta and Yuming Zhao. The record does not separately identify a proposer.
\item Supplied literature source: Mehta--Slofstra--Zhao, ``positivity is undecidable in tensor products.'' \url{https://doi.org/10.4171/JEMS/1786}. This is the supplied reference label; a complete bibliographic citation and precise theorem statement are not provided.
\item Supplied literature source: ``OWR formulation and free-group special case.'' \url{https://doi.org/10.4171/owr/2021/49}.
\item Literature assessment checked at \texttt{2026-08-22}, exactly as recorded. This is supplied dated curation, not independently verified. It reports no single test for arbitrary star-algebras, effective procedures for free-group positivity, and coRE-hardness and undecidability for positivity in tensor products of free algebras. The supplied background expressly requests independent verification of time-sensitive status and cited claims.
\item The registry labels the statement \texttt{corrected\_verified} and says the official report was checked directly and the cleaned statement preserves the complete mathematical target. Those are supplied verification claims; no independent source consultation is asserted in this draft.
\end{itemize}
\end{document}
